\subsection{事件—区域—覆盖：最小时空骨架先于几何}\label{subsec:physical-spacetime-skeleton-event-region-cover}

如果说上一小节把观察者从主体改写为状态纤维，那么接下来必须完成的工作，就是把空间从点集坐标改写为区域—覆盖系统。因为一旦继续用“点上发生了什么”来理解全文后续的对象、值、投影、谱与资源结构，就会把真正有逻辑意义的局部性掩盖掉。第 \ref{sec:logic-expansion-chain} 节已经表明，空间 forcing 的基本单位不是点，而是区域；区域之间的最重要关系不是距离，而是包含、覆盖、限制与拼接。因此，本章在物理化地重读这些结构时，也必须坚持这个顺序：事件—区域—覆盖的最小时空骨架先于几何坐标。

这里所谓事件，并不是指一个带坐标的点状事实，而是指一个在某个区域上已经进入对象层决定包络的局部断言。换言之，一个事件之所以成立，不在于它占据了某个绝对位置，而在于它在某个局部区域中已经具备了可比较、可拼接、可传递的结构条件。从这个角度看，第 \ref{sec:recursive-addressing} 节中的地址对象、局部截面与 \(\mathsf{NULL}\)，实际上就是最原始的事件层：一个地址若尚未对象化，它连事件资格都没有；一个局部截面若不存在，则事件在该区域上仍然是空的；只有当一个对象在某个区域上具有合法、可核验的局部截面时，它才真正成为一个“发生了”的局部事件。也就是说，事件性并不先验属于世界，而是由局部 forcing 赋予。

一旦这样理解，区域与覆盖的重要性就会完全显现。区域不是某个坐标盒子，而是“事件能够在其中被局部表达并与邻域拼接”的最小承载体。覆盖则不是一个几何附属结构，而是说明整体事件如何由若干局部事件兼容地拼出来的规则。第 \ref{sec:logic-expansion-chain} 节中的层结构里的拼接公理、第 \ref{sec:recursive-addressing} 节里的前缀站点与局部截面、第 \ref{sec:pom} 节里的投影词接口，全部都可在这一层重新读成同一件事：整体不是先验整块给定，而是由兼容局部在覆盖之下被唯一重建出来的。于是，本章中的时空骨架也不应从点—线—面的几何对象出发，而应从局部事件—区域支撑—覆盖拼接的逻辑对象出发。

从这里进一步可以得到一个非常强的结论：几何在本章不是第一性的，而是第二性的。若某个后续章节真的引入了具体距离、相位、半径、窗口或投影门，那么这些量都应被理解为对已有区域—覆盖骨架的某种精细化或参数化，而不是对空间 forcing 的替代。也就是说，真正的最小时空骨架，不是一个现成的流形，而是一套允许局部事件进入对象层、允许局部证书被比较、允许整体由覆盖拼接恢复的结构。只有在这一基础上，后文关于相位门、谱接口、primitive 轨道、资源函数与物理化读法才不会失去其逻辑根基。换言之，本小节的核心命题应写成：空间首先是事件可以被局部表达、限制与拼接的区域系统；几何只是这一系统的进一步参数化，而不是其起点。
